\paragraph{Problem representation.}
The simulator does not choose arbitrary activity labels at a decision point.
The BPMN engine first converts the current case marking into executable
\emph{transition candidates}. Each candidate contains the BPMN transition id,
the visible activity label and the silent-transition path used to reach the
visible transition. Branching is therefore constrained by the process model:
every returned decision is filtered against the currently enabled BPMN
candidates before it can be fired.

\paragraph{Composite decision chain.}
The final branching chain used in the evaluated simulator is
\[
    \text{Predictive} \rightarrow \text{Probability} \rightarrow \text{Random}.
\]
The corrected repository separates the leakage-free evaluation artifact from
the full-log deployment artifact. The final optimized evaluation candidate is
\texttt{composite\_branching\_evaluation\_train70\_rfopt\_v1.pkl}; the matching
full-log artifact is retained only for deployment. Held-out claims and the
candidate fixed replay use the evaluation-scoped artifact.
This chain is the runtime policy, not the unit of the offline classifier score:
the offline tables evaluate the predictive model separately, while the simulator
uses the composite chain to remain robust whenever the ML prediction cannot be
mapped to a currently enabled BPMN candidate.
The predictive component is a Random-Forest next-activity classifier, following
the ensemble-learning idea of Breiman~\cite{breiman2001randomforests} and the
scikit-learn implementation interface~\cite{pedregosa2011scikit}. Its offline
training examples are obtained by replaying the event log on
\texttt{models/v4\_replay.bpmn}. A supervised row is emitted only after the
observed current event has been synchronized and fired in the BPMN marking, the
marking exposes multiple BPMN-valid next labels, and the observed next event is
one of those enabled labels. The RF-optimized feature vector combines
control-flow history (current and previous activity and prefix information),
time/resource context and non-constant case attributes such as application
type, loan goal and requested amount. At runtime, the classifier
ranks possible next activities, but the selected label is accepted only if it is
BPMN-valid for the current marking. The RF-optimized evaluation artifact keeps
the non-constant runtime-equivalent features and drops constant counters from
the candidate feature set. If no valid predictive label remains, the engine
calls the probabilistic fallback. Therefore, the implementation satisfies the
replay-based interpretation of the advanced branching requirement for the
subset of event-log behavior that can be synchronized with the provided BPMN
model; the remaining coverage is reported separately.

The probabilistic component estimates empirical successor distributions from
training data only. It first tries a state-conditioned distribution keyed by
current activity, previous activity and coarse repetition counters. If that
state has insufficient support, it falls back to the global distribution for the
current activity and finally to a candidate-uniform BPMN-valid fallback. The
final contribution therefore consists of a clean probabilistic baseline, a
history-aware Random-Forest predictor, and a composite runtime chain that keeps
all choices BPMN-valid.

\paragraph{BPMN alignment.}
The transition-aware integration separates event-log labels from BPMN transition
ids. This matters because a label can be ambiguous or temporarily non-executable
under the current marking. The BPMN engine exposes transition candidates rather
than labels alone, and the integration records whether a historical or predicted
activity can be synchronized with an identifiable transition. Consequently,
next-label prediction, label-to-transition coverage and exact transition fire
success are reported as separate quantities in Sec.~\ref{ssub:branching-eval};
the latter is not treated as transition-classification accuracy because the log
does not provide transition-ID ground truth.
